\section{Related literature}\label{sec:lit}

Since the Bayh--Dole Act placed ownership of federally funded inventions with universities, the question of how university
rules shape commercialization has been central to the study of academic technology transfer
\citep{MoweryEtAl2001,HendersonJaffeTrajtenberg1998,Sampat2006,GrimaldiEtAl2011}. Institutional ownership created both the
authority and the obligation to manage inventions, and universities built policies and offices to exercise it. A first
body of evidence asks whether the terms of those policies matter. Universities with higher inventor royalty shares earn more
licensing income \citep{LachSchankerman2008}; willingness to take equity and lower royalty shares are associated with more
startups \citep{DiGregorioShane2003}; and public universities with local-development objectives license differently from
private ones \citep{BelenzonSchankerman2009}. Related evidence from Europe links royalty sharing to inventive effort
\citep{ArqueCastellsEtAl2016}.

These studies compare institutions, and institutions that adopt particular terms differ in many other ways. When policy
data are corrected and expanded, the royalty-share relationship weakens considerably \citep{OuelletteTutt2020}, which
suggests that single contractual terms observed across universities are fragile guides to what policy changes would do. The
most convincing evidence therefore comes from reforms that change rules within a system. \citet{HvideJones2018} show that
when Norway transferred ownership of university inventions from researchers to institutions, startups and patenting by
university researchers fell by roughly half, and \citet{GeunaRossi2011} document the broader European shift toward
institutional ownership. \citet{KenneyPatton2009} argue on similar grounds that alternatives to the prevailing ownership
model could improve commercialization. Ownership reforms, however, are rare and large. They change who holds the rights to
an invention and therefore operate first on researchers' incentives to commercialize.

Most policy change that universities undertake is of a different kind. Institutions periodically rewrite their IP policies,
usually retaining institutional ownership while changing disclosure procedures, revenue sharing, conflict-of-interest rules,
equity provisions, and the language through which the relationship between inventors and the institution is described.
These ordinary revisions are the policy instrument that universities and state systems actually control, yet they have not
been studied as events. Whether they are followed by changes in commercialization, and at which stage, is an open question.

Where such revisions might matter depends on how commercialization is organized. A large literature treats the TTO as the
locus of commercialization performance, with productivity that varies with staffing, incentive pay, routines, and
relationships with industry \citep{SiegelWaldmanLink2003,SiegelWaldmanAtwater2004,ThursbyKemp2002,MarkmanEtAl2005,
OsheaEtAl2005}. TTOs are boundary-spanning organizations that translate between science, law, and markets
\citep{DebackereVeugelers2005,PerkmannEtAl2013,HolgerssonEtAl2019}, and most university inventions are licensed at an
embryonic stage, so licensing depends on continued cooperation among inventors, the office, and licensees
\citep{JensenThursby2001}. Formal policy is one of the instruments through which this cooperation is governed.

This organizational view distinguishes two margins at which policy change could operate. The first is faculty entry into the
formal process. Many potentially commercial inventions are never disclosed \citep{JensenThursbyThursby2003}, disclosure
depends on local norms, peers, and perceptions of the TTO \citep{OwenSmithPowell2001,BercovitzFeldman2008}, and much of the
growth in university licensing has been attributed to faculty's increasing willingness to disclose
\citep{ThursbyThursby2002}. If a revision changed how faculty perceive the process, disclosures should respond first. The
second margin lies after disclosure, where the office evaluates, protects, markets, and negotiates. A revision that changed
how the TTO and inventors interact with licensees would appear in licensing without a prior change in disclosure. Stage-specific
measures make it possible to tell these apart; the filing ratio of new patent applications to disclosures, for example,
isolates the decision to seek protection for what has been disclosed \citep{ChoudhryPonzio2020,Hall2021}.

Whether written policy affects either margin is itself contested. Institutional theory holds that formal structures can be
ceremonial, adopted for legitimacy and decoupled from practice \citep{MeyerRowan1977}. Research on administrative burden
argues the opposite: the learning, compliance, and psychological costs embedded in formal procedures shape how people
interact with institutions \citep{MoynihanHerdHarvey2015,HerdMoynihan2018}. Policy text is also a record of organizational
choices that may travel with other changes, such as new office leadership or reorganization. The relationship between
rewriting a policy and subsequent behavior is therefore an empirical question.

Answering it requires two things the literature has lacked: dated revisions of the policies universities actually use, and
a consistent characterization of what each revision changed. The companion paper provides the second by measuring the
communicative stance of policy text \citep{SharmaPCSI2026}, building on methods for measuring institutional documents as
text \citep{LoughranMcDonald2011,GentzkowKellyTaddy2019,GrimmerRobertsStewart2022}. This paper supplies the first and asks
what follows. Relative to cross-institution studies, it compares each university with itself before and after a revision
and with contemporaneous non-revisers. Relative to ownership reforms, it studies the ordinary revisions that universities
choose. Relative to the TTO literature, it asks whether formal policy change is followed by change at the office's
licensing margin or at faculty disclosure, and so connects formal rules to the organizational account of technology
transfer.
